\section{Discussion}
\label{sec:discussion}

\subsection{Why certification improves reliability and when it adds no cost}
The experiments show why point-estimate tuning fails:
selecting the cheapest configuration whose measured risk clears a
threshold is an uncorrected multiple-testing procedure, so its winner
is systematically the configuration with favorable measurement noise,
an instance of the winner's curse. The effect is strongest where adaptive
systems are most attractive: many candidates, small quality gaps
between adjacent tiers, and moderate audit budgets. Certification
neutralizes the curse not by testing less but by testing in a
data-independent order and stopping at the first failure, which caps
the family-wise error at $\delta$ regardless of how many candidates
follow. The cost of the guarantee is a finite-sample margin of
0.02--0.07 at practical calibration sizes --- and on benchmarks where
the certified family is richer than the baselines' (CRAG), the
certified policy is simultaneously the safest and among the cheapest,
so certification adds no cost in those settings. The margin shrinks at the standard
$O(n^{-1/2})$ rate, so operators can trade audit budget for cost
directly.

\subsection{Learning mechanism and complementarity}
The proposed separation between learning and certification is useful
for adaptive inference. Sufficiency models compress retrieval margins,
support, citation, age, and abstention signals into a stopping score.
Better scores move more easy queries to early rungs and improve cost.
However, Theorem~\ref{thm:validity} does not require the score to be
calibrated or even informative. This is deliberate: prediction quality
determines efficiency, while the held-out test determines validity.
The score ablation makes the distinction visible. Removing informative
signals reduces the savings and shortens the certified prefix, but it
does not increase the allowed certification error.

ContractRAG is therefore not an alternative to RAG-LER, IRAGKR,
AIR-RAG, or uncertainty-aware retrieval. Those methods learn stronger
rerankers, query transformations, and retrieval gates. They can be used
as rungs or features in the policy family, improving its attainable
frontier. ContractRAG contributes the selection rule above that family.
The dense-grid and policy-family studies show that a richer frontier is
useful only when its ordering retains enough power; validity itself is
unchanged. Finally, a learned gate is tied to the distribution on which
it was fitted. The restart-mixture detector supplies the missing
temporal mechanism by deciding when the old certificate and score
should be replaced with a fresh epoch.

\subsection{Limitations and future directions}
The study has several limitations. (i) The guarantee is relative to
the audited loss: a systematically biased auditor (an LLM judge that
overrates fluent answers) shifts what is being certified; the
procedure also accepts human audits, and the audit
accounting of Section~\ref{sec:exp-drift} shows their cost is a
first-order term that must be (and in our tables is) charged against
the savings. (ii) Certification assumes the calibration sample is
exchangeable with deployment within an epoch; adversarial or rapidly
cycling drift would trigger frequent recalibration, and although the
geometric spending schedule keeps lifetime validity, each cycle halves
the remaining budget and lengthens detection delays logarithmically
--- safe mode between an alarm and the next successful recertification
is best-effort, not certified. (iii) The fixed-sequence procedure
returns the cheapest policy of its certified prefix; this coincides
with the family-wide cheapest feasible policy only under the ordering
conditions of Proposition~\ref{prop:consistency}, and the dense-grid
experiment shows the power cost of violating them. (iv) The policy
families we certify are built over a linear ladder plus routers;
richer spaces (per-operator adaptive plans, multi-agent pipelines)
enlarge the search but also stress the power of a fixed-sequence walk
at small audit budgets. (v) Group-conditional certification requires
groups known at query time and pays a Bonferroni factor; overlapping
or learned groups are future work. (vi) Our cost model is monetary;
extending contracts to energy or carbon budgets is straightforward but
unexplored.

\subsection{Threats to validity}
\emph{Internal validity}: all methods consume identical execution
matrices, so comparisons cannot be confounded by retrieval variance;
the repeated-draw protocol (1000 draws) evaluates each selected policy
on the \emph{exact} finite population, so the reported violation
probabilities carry Monte Carlo error of about $\pm 0.014$ and no
test-sampling error at all. \emph{Statistical validity}: the protocol
samples calibration sets without replacement while the certificate's
p-value assumes i.i.d.\ draws. As discussed in Section~\ref{sec:setup},
the p-value remains conservative in the finite-population regime, and
the exact hypergeometric refinement only strengthens the results.
\emph{Construct validity}: losses follow each benchmark's official
protocol (token-F1, CRAG's LLM-judge grading, ALCE's TRUE-NLI citation
entailment), whose agreement with human raters is documented by their
authors; the certificate is exact for these losses, and
Section~\ref{sec:discussion} discusses auditor bias. \emph{External
validity}: results cover four benchmarks, four LLM families from four
model developers, two serving backends (one hosted API platform, one
local deployment), and two cost currencies, but all workloads are
English-language QA; agentic and multilingual workloads may distribute
difficulty differently. \emph{Reproducibility}: every LLM call is
cached token-exactly, and all numbers regenerate from the released
artifact without further API calls.
